% ==============================================================================
% 公共题库: common/pool/解三角形/难题/q_triangle_dof_vector_median.tex
% 难度: ★★★ 难题
% 知识点: 中线倍长、平行四边形恒等式、自由度传递、向量数量积
% ==============================================================================

\newcommand{\qTriangleDofVectorMedianStem}{%
在 \(\triangle ABC\) 中，\(b = 4\)，\(c = 3\)，\(AD\) 是 \(BC\) 上的中线，已知 \(AD = m_a = \frac{\sqrt{19}}{2}\)。
\begin{enumerate}[(1)]
  \item 求边长 \(a\) 的值；
  \item 求 \(\cos A\) 的值；
  \item 若点 \(G\) 为 \(\triangle ABC\) 的重心，求向量数量积 \(\vec{GA} \cdot \vec{GB}\) 的值。
\end{enumerate}%
}

\newcommand{\qTriangleDofVectorMedianAnswer}{%
(1) \(a = 5\)；\qquad (2) \(\cos A = 0\)（\(A = 90^\circ\)）；\qquad (3) \(\vec{GA} \cdot \vec{GB} = -\frac{2}{9}\)。
}

\newcommand{\qTriangleDofVectorMedianSolution}{%
(1) 倍长中线得平行四边形恒等式 \(2(b^2+c^2) = a^2 + (2m_a)^2 \implies 2(25) = a^2 + 19 \implies a=5\)。
(2) 由余弦定理 \(\cos A = \frac{4^2+3^2-5^2}{2 \cdot 4 \cdot 3} = 0 \implies A = 90^\circ\)。
(3) 重心向量分解：\(\vec{GA} = -\frac{1}{3}(\vec{AB}+\vec{AC})\), \(\vec{GB} = \frac{2}{3}\vec{AB}-\frac{1}{3}\vec{AC}\)。
\(\vec{GA} \cdot \vec{GB} = -\frac{2}{9}|\vec{AB}|^2 + \frac{1}{9}|\vec{AC}|^2 = -\frac{2}{9}(9) + \frac{1}{9}(16) = -\frac{2}{9}\)。
}

\newcommand{\qTriangleDofVectorMedianDiagnosis}{%
【错因】未掌握倍长中线公式，对重心向量与主向量代数分解不够熟练。
}
